
\lesson{19/10}
\textbf{Note} Translations are always transitive, i.e. there is only
one orbit which consists of the entire group, since
$$g\cdot x=y\implies g = y\cdot x^{-1}$$
or, equivalently,
$$\forall x,y\in G,\,\exists g\in G\text{ s.t. }L_g(x)=y$$
\textbf{Proof} Conjugation is a left action.\\
Consider $L_{g\cdot g'}$:
$$L_{g\cdot g'}(x) = (g\cdot g')x(g\cdot g')^{-1} = 
g\cdot(g'\cdot x \cdot (g')^{-1})\cdot g^{-1} = L_g(L_{g'}(x)) = L_g\circ L_{g'}
(x)\qed$$

Conjugation is transitive if $G$ is Abelian, but not in general.

\textbf{Def.} Given $g\in G$, the translation action $L$ and a proper
subgroup $H\subset G$, the set
$$gH := \left\{g'\in G\text{ s.t. }\exists h\in H\text{ s.t. }g'=L_g(h)=g\cdot h
\right\}$$
is called a \textit{left coset of} $H$.

If $H$ is a proper subgroup, there exist more than one left coset. The set of
all left cosets is called the \textit{quotient space}:
$$G/H = \left\{gH,\,\forall g\in G\right\}$$
\textbf{Note} This terminology is akin to that of equivalence relations, because
the equivalence relation "$g$ is in relation to $g'$ if, for some $h\in H$,
$g'=g\cdot h$" forms equivalence classes $gH$ (with representative $g$) on $G$,
which are disjoint and cover all of $G$.

\textbf{Proposition} \textit{The elements of $gH$ are in one-to-one correspondence
with those of $H$, thus $|H|=|gH|$}.\\
\textbf{Proof} Consider $g\in G$ and the function
\begin{align*}
    f_g:&H\to gH\\
        &h\to g\cdot h:
\end{align*}
$f_g$ is injective, since
$$f_g(h_1)=f_g(h_2) \implies g\cdot h_1 = g\cdot h_2 \iff h_1 = h_2$$
and surjective (by definition of $gH$): so $f_g$ is a one-to-one correspondence.
\qed

\textbf{Lagrange theorem} \textit{The order $|H|$ of any proper subgroup
$H\subset G$ must be a divisor of $|G|$}.\\
\textbf{Proof} Since $gH$ are equivalence classes of $G$, they are disjoint, so
$$|G| = \sum_{g}^{}|gH|:$$
for the previous proposition, $\forall g,g'\in G,\,|gH|=|H|=|g'H|$ and so
$$|G|=n|gH| = n|H|,\;\text{with }n\in\mathbb{N}_0 \qed$$
